\documentclass[11pt,oneside,openany]{book}

\usepackage[T1]{fontenc}
\usepackage{lmodern}
\usepackage[letterpaper,margin=1.10in,includeheadfoot]{geometry}
\usepackage{amsmath,amssymb,amsthm,mathtools}
\usepackage{microtype}
\usepackage{xcolor}
\usepackage{titlesec}
\usepackage{fancyhdr}
\usepackage{needspace}
\usepackage{hyperref}

\ifdefined\pdftrailerid
  \pdftrailerid{MathGenome ALG-01 core supplement}
\fi

\definecolor{ink}{HTML}{172923}
\definecolor{deepgreen}{HTML}{2B5242}
\definecolor{rulegreen}{HTML}{8A9F92}
\definecolor{palegreen}{HTML}{F0F4EF}
\hypersetup{
  colorlinks=true,linkcolor=deepgreen,urlcolor=deepgreen,citecolor=deepgreen,
  pdftitle={Linear and Multilinear Algebra and Matrix Theory},
  pdfsubject={Symmetric powers over commutative rings and Witt theory over fields},
  pdfauthor={MathGenome},bookmarksnumbered=true
}
\setcounter{tocdepth}{1}
\setcounter{secnumdepth}{2}
\setlength{\headheight}{15pt}
\setlength{\emergencystretch}{2em}
\titleformat{\chapter}[display]
  {\normalfont\color{ink}}{\color{deepgreen}\large\scshape Chapter \thechapter}{9pt}
  {\huge\bfseries}
\titlespacing*{\chapter}{0pt}{4pt}{24pt}
\titleformat{\section}{\normalfont\Large\bfseries\color{ink}}{\thesection}{0.7em}{}
\titleformat{\subsection}{\normalfont\large\bfseries\color{ink}}{\thesubsection}{0.7em}{}

\pagestyle{fancy}
\fancyhf{}
\fancyhead[L]{\small\sffamily\color{deepgreen}MathGenome}
\fancyhead[R]{\small\nouppercase{\leftmark}}
\fancyfoot[C]{\small\thepage}
\renewcommand{\headrulewidth}{0.35pt}
\renewcommand{\chaptermark}[1]{\markboth{\thechapter.\ #1}{}}
\fancypagestyle{plain}{%
  \fancyhf{}\fancyfoot[C]{\small\thepage}\renewcommand{\headrulewidth}{0pt}}

\newtheorem{theorem}{Theorem}[chapter]
\newtheorem{lemma}[theorem]{Lemma}
\newtheorem{proposition}[theorem]{Proposition}
\newtheorem{corollary}[theorem]{Corollary}
\theoremstyle{definition}
\newtheorem{definition}[theorem]{Definition}
\newtheorem{example}[theorem]{Example}
\theoremstyle{remark}
\newtheorem{remark}[theorem]{Remark}
\allowdisplaybreaks[1]

\begin{document}
\frontmatter
\hypersetup{pageanchor=false}
\begin{titlepage}
  \thispagestyle{empty}
  \thispdfpagelabel{Cover}
  \begin{flushleft}
    \vspace*{0.65in}
    {\large\scshape\color{deepgreen}MathGenome\par}
    \vspace{0.32in}
    {\color{rulegreen}\rule{0.94\textwidth}{0.65pt}\par}
    \vspace{0.45in}
    {\color{ink}\fontsize{34}{39}\selectfont\bfseries
       Linear and Multilinear\par Algebra and Matrix Theory\par}
    \vspace{0.30in}
    {\color{deepgreen}\Large
       Symmetric Powers and Witt Theory\par}
    \vfill
    {\large Supplementary lecture notes\par}
    \vspace{0.12in}
    {\normalsize ALG-01 \enspace\textperiodcentered\enspace MathGenome v0.22.0 reference baseline\par}
    \vspace*{0.48in}
  \end{flushleft}
\end{titlepage}
\hypersetup{pageanchor=true}
{\small\tableofcontents}
\chapter*{Scope and reading guide}
\addcontentsline{toc}{chapter}{Scope and reading guide}
\markboth{Scope and reading guide}{}

This supplement develops symmetric powers over arbitrary commutative
rings and Witt decomposition and cancellation over fields of
characteristic different from two. The route through ALG-01 below
points to the existing proofs and identifies where the new chapters
enter.

{\setlength{\fboxsep}{0.8em}\setlength{\fboxrule}{0.35pt}%
\noindent\fcolorbox{rulegreen}{palegreen}{%
\begin{minipage}{\dimexpr\linewidth-2\fboxsep-2\fboxrule\relax}
\textbf{Result keys.} \texttt{LA01-C05.01} belongs to module
\texttt{ALG-01-C05}. The source map records its conventions,
dependencies, and PDF and TeX locations. Keys beginning with
\texttt{K05} refer to result groups in the released D05 notes.
\end{minipage}}\par}

\section*{The six core modules}

\paragraph{C01 \enspace Vector spaces, modules, and dimension.}
D05, groups K05.02--03 and K05.06, proves the field results on bases,
rank--nullity, and quotients \cite{d05}.
\emph{Algebra Core} supplies module quotients, exactness, and split
sequences \cite[Ch.~1]{p103}. D07 treats modules over a PID \cite[Ch.~16]{d07}.

\paragraph{C02 \enspace Duality and forms.}
D05, groups K05.07 and K05.11, gives dual bases, the dual of a quotient,
scalar conventions on forms, and nondegeneracy.
Group K05.13 gives the adjoint on finite-dimensional real and complex
inner-product spaces.

\paragraph{C03 \enspace Determinants and canonical forms.}
D05, groups K05.05, K05.08, and K05.10, gives determinants, Cayley--Hamilton,
primary decomposition, and rational and Jordan forms.
Jordan form uses a splitting field. For the Smith form over a PID,
see D07, Chapter~16.

\paragraph{C04 \enspace Spectral and singular-value methods.}
D05, groups K05.12--14, proves the finite-dimensional real and complex
spectral, SVD, and polar results. For the Hilbert-space partial-isometry
version see \cite{d16}, and for the min--max and matrix perturbation
bounds see \cite{p206}.

\paragraph{C05 \enspace Tensor, exterior, and symmetric powers.}
The balanced tensor construction is in \cite[Ch.~1]{p103}. D05,
group K05.15, treats the field case. The comparison
\(S\otimes_R\bigwedge_R^n M\simeq\bigwedge_S^n(S\otimes_R M)\)
is the degree-\(n\) comparison in \cite[RK-03 and RK-04]{commalg}.
Chapter~\ref{ch:la01-symmetric} develops the symmetric powers over rings.

\paragraph{C06 \enspace Quadratic forms and Witt theory.}
D05, group K05.11, provides polarization, diagonalization, and the
orthogonal-complement results. Chapter~\ref{ch:la01-witt} gives
hyperbolic splitting and Witt cancellation with the characteristic
hypothesis and the nondegenerate scope planned for NT-06.
For the classical treatment see \cite[Ch.~I, Sec.~4]{lam-witt}.

\mainmatter
\chapter{Symmetric powers over a commutative ring}\label{ch:la01-symmetric}

Throughout the chapter \(R\) is a commutative ring with identity, all
algebra maps preserve the identity, and \(M\) is an \(R\)-module.
The construction uses the balanced tensor product and its
associativity from \emph{Algebra Core}
\cite[ALG-03-C02-S01]{p103}. The field case, including the monomial
basis, is developed in \emph{Linear Algebra}
\cite[entries 15.symmetric and 15.symmetric\_basis]{d05}. Here the
coefficient ring and the module may be arbitrary.

\section{The symmetric algebra and its universal properties}

\begin{definition}[LA01-C05.01. The symmetric algebra]
\label{def:la01-symmetric-algebra}
The tensor algebra is \(T_R(M)=\bigoplus_{n\ge 0}M^{\otimes_R n}\),
with \(M^{\otimes_R 0}=R\) and product given by concatenation of
tensors. The symmetric algebra \(\operatorname{Sym}_R(M)\) is its
quotient by the two-sided ideal \(J\) generated by
\(x\otimes y-y\otimes x\) for \(x,y\in M\). The degree-\(n\) component
is the symmetric power \(\operatorname{Sym}^n_R(M)\).
\end{definition}

The ideal \(J\) is homogeneous because its generators have degree
two. More explicitly, a homogeneous piece of a product
\(a(x\otimes y-y\otimes x)b\) is a sum of such products where \(a,b\)
are homogeneous and their degrees add to the required degree minus
two. There are no relations in degrees zero and one. Thus
\(\operatorname{Sym}^0_R(M)=R\) and
\(\operatorname{Sym}^1_R(M)=M\) canonically. We write \(x_1\cdots x_n\)
for the image of a pure tensor, with empty product \(1\). Scalars
commute with everything, and degree-one generators commute by the
defining relations. Since these generate the quotient as an
\(R\)-algebra, it is commutative. See \cite[Algebra, Section 10.13]{stacks-sym}.

\begin{proposition}[LA01-C05.02. Universal symmetric multilinear maps]
\label{prop:la01-symmetric-multilinear}
For every \(n\ge 0\) and every \(R\)-module \(N\), the assignment
\(g\mapsto ((x_1,\ldots,x_n)\mapsto g(x_1\cdots x_n))\) gives a natural
bijection from \(R\)-linear maps \(\operatorname{Sym}^n_R(M)\to N\)
to symmetric \(R\)-multilinear maps \(M^n\to N\). For \(n=0\) the
right-hand side means the choices of a single element of \(N\).
\end{proposition}
\begin{proof}
For \(n\ge 1\), iteration of the tensor universal property
gives a unique linear map on \(M^{\otimes_R n}\) extending any
multilinear map on \(M^n\). In degree \(n\ge 2\), the homogeneous
description of \(J\) shows that \(J\cap M^{\otimes_R n}\) is spanned
by differences of pure tensors related by one adjacent swap.
Indeed the factors before and after the generator can be taken
to be pure tensors, which span their respective tensor powers.
A linear map kills these differences precisely when its associated
multilinear map is invariant under adjacent swaps. Those swaps
generate the whole permutation group. This proves the bijection
for \(n\ge 2\), and for \(n=1\) both descriptions give all linear
maps \(M\to N\). For \(n=0\) use \(g\mapsto g(1)\) on
\(\operatorname{Hom}_R(R,N)\).

The construction is natural in \(N\) by postcomposition. A
linear map \(M\to M'\) induces tensor powers that take adjacent
differences to adjacent differences. The induced maps on the
quotients give naturality in \(M\) by precomposition.
\end{proof}

\begin{proposition}[LA01-C05.03. The free commutative algebra property]
\label{prop:la01-symmetric-free}
For any commutative \(R\)-algebra \(A\), each \(R\)-linear map
\(f\colon M\to A\) extends uniquely to an \(R\)-algebra map
\(\operatorname{Sym}_R(M)\to A\).
\end{proposition}
\begin{proof}
For every \(n\), the product \(f(x_1)\cdots f(x_n)\) is symmetric
and multilinear. Proposition~\ref{prop:la01-symmetric-multilinear}
extends it to a linear map on the degree-\(n\) component.
For degree zero this gives the structure map \(R\to A\).
The direct sum of these maps respects multiplication on products
of pure classes, where it agrees with multiplication of the images
of all their factors in \(A\). By bilinearity it respects
multiplication on the whole algebra, and its degree-zero component
takes the unit to the unit. The images of \(R\) and of \(M\)
determine the map on all products and linear combinations of
generators, giving uniqueness.
\end{proof}

\begin{theorem}[LA01-C05.04. The monomial basis of a finite free module]
\label{thm:la01-symmetric-basis}
Suppose \(M\) is free with basis \(e_1,\ldots,e_r\), allowing \(r=0\).
The assignment \(e_i\mapsto X_i\) gives a graded \(R\)-algebra
isomorphism
\[
  \operatorname{Sym}_R(M)\cong R[X_1,\ldots,X_r],
\]
where every \(X_i\) has degree one. For every \(n\ge 0\), the
elements \(e_1^{a_1}\cdots e_r^{a_r}\), with all \(a_i\ge 0\) and
\(\sum_i a_i=n\), are an \(R\)-basis of \(\operatorname{Sym}^n_R(M)\).
When \(r\ge 1\) there are \(\binom{n+r-1}{n}\) basis indices. For
\(r=0\), the basis consists of \(1\) in degree zero and is empty
in positive degrees.
\end{theorem}
\begin{proof}
Send the basis linearly to the variables and extend using
Proposition~\ref{prop:la01-symmetric-free}. Sending \(X_i\) to \(e_i\)
defines the reverse map by the universal property of the polynomial
algebra. Each composition fixes the scalars and all the algebra
generators on its side, so both compositions are identities.
The maps preserve degree. The formal polynomial algebra is the free
\(R\)-module on its formal monomials, and its degree pieces are
spanned by the monomials of the indicated total degree. For \(r\ge 1\),
the exponents can be recorded by placing \(r-1\) separators among
\(n+r-1\) positions, with the \(n\) other positions recording the
degree. For \(r=0\) the polynomial algebra is \(R\) in degree zero.
This gives the asserted bases in every degree.
\end{proof}

The basis statement over a commutative ring also follows from the
multilinear approach in \cite[Theorem 1.2]{conrad-sym}.

\section{Extension of scalars}

\begin{theorem}[LA01-C05.05. Arbitrary base change]
\label{thm:la01-symmetric-base-change}
Let \(R\to S\) be any homomorphism of commutative rings with
identity and let \(M_S=S\otimes_R M\). The graded \(S\)-algebra map
\[
  S\otimes_R \operatorname{Sym}_R(M)
  \longrightarrow \operatorname{Sym}_S(M_S),
  \qquad
  s\otimes (x_1\cdots x_n)\longmapsto
       s(1\otimes x_1)\cdots(1\otimes x_n)
\]
is an isomorphism and gives a comparison in every degree \(n\ge 0\).
It is natural in \(M\) and compatible with compositions of
ring maps under the usual associativity of scalar extension.
\end{theorem}
\begin{proof}
The \(R\)-linear map \(x\mapsto 1\otimes x\), with values in
degree one of \(\operatorname{Sym}_S(M_S)\), extends to an
\(R\)-algebra map by Proposition~\ref{prop:la01-symmetric-free}.
Multiplying its values by \(s\) defines the indicated map on
simple tensors, by the balanced tensor universal property. This
map respects scalar multiplication and algebra multiplication.
In the opposite direction, \(s\otimes x\mapsto s\otimes x\) with
values in degree one of \(S\otimes_R\operatorname{Sym}_R(M)\)
is the scalar extension of the degree-one inclusion. It extends
to an \(S\)-algebra map by the same proposition over \(S\).
The \(1\otimes x\) generate both algebras over \(S\), and the
two compositions fix these elements and the scalars. The maps
are therefore inverse. They preserve degree, and tensor products
commute with direct sums, so the degreewise comparisons follow.

For a module map \(f\colon M\to M'\), the two routes send the
generator \(1\otimes x\) to \(1\otimes f(x)\). For \(S\to U\),
the two routes after scalar extension send \(1\otimes(1\otimes
x)\) to the degree-one generator \(1\otimes x\) over \(U\).
In both cases agreement on the generators proves the claim.
\end{proof}

The base-change comparison is the module version of the pullback
formula for symmetric algebras
\cite[Modules, Lemma 17.21.3]{stacks-sym}. For exterior powers,
\emph{Advanced Commutative Algebra} gives the alternating quotient
construction and its coefficient comparison in RK-03 and RK-04
\cite[RK-03 and RK-04]{commalg}. In degree \(q\) the map is
\[
  s\otimes(x_1\wedge\cdots\wedge x_q)
  \longmapsto s(1\otimes x_1)\wedge\cdots\wedge(1\otimes x_q).
\]
The alternating condition is vanishing on repeated inputs,
including in characteristic two. RK-04 applies the comparison to
the terms of the Koszul complex for every ring map and uses
flatness in the statement about Koszul homology.

\begin{example}[LA01-C05.06. A torsion module and nonflat base change]
\label{ex:la01-symmetric-torsion}
Take \(R=\mathbb Z\), \(M=(\mathbb Z/2\mathbb Z)e\) with generator
\(e\), and \(S=\mathbb F_2\). A linear map \(M\to A\) into a
commutative \(\mathbb Z\)-algebra is specified by an element of
\(A\) annihilated by \(2\). The same choices specify
\(\mathbb Z\)-algebra maps \(\mathbb Z[X]/(2X)\to A\) via \(X\).
The canonical maps induced by \(e\leftrightarrow X\) are inverse
on the generators, so
\(\operatorname{Sym}_{\mathbb Z}(M)\cong\mathbb Z[X]/(2X)\).
Its degree-zero part is \(\mathbb Z\), and each positive degree
is generated by \(X^n\) with coefficient module \(\mathbb Z/2\mathbb Z\).
After scalar extension the quotient becomes \(\mathbb F_2[X]\),
in agreement with the rank-one description of
\(\operatorname{Sym}_{\mathbb F_2}(M_S)\).
This extension is nonflat because tensoring
\(\mathbb Z\xrightarrow{\,2\,}\mathbb Z\) with \(\mathbb F_2\)
sends an injection to the zero map on a nonzero module.
\end{example}

\begin{example}[LA01-C05.07. The quotient in characteristic two]
\label{ex:la01-symmetric-char-two}
Let \(M=\mathbb F_2 e_1\oplus\mathbb F_2 e_2\) and let \(\tau\)
exchange the factors of \(M^{\otimes 2}\) over \(\mathbb F_2\).
In the tensor basis the fixed subspace has basis
\[
  e_1\otimes e_1,\qquad e_2\otimes e_2,\qquad
  e_1\otimes e_2+e_2\otimes e_1.
\]
Indeed the coefficients of the two mixed tensors must be equal.
The quotient map takes the third basis element to
\(2e_1e_2=0\), even though \(e_1e_2\) is a nonzero basis monomial
of the quotient. Thus the map from the fixed tensors to the
symmetric quotient is neither injective nor surjective. The
averaging construction in \cite[entry 15.normalization]{d05}
uses invertible factorials in its field setting, while the
quotient and the monomial basis above hold in every characteristic.
\end{example}
\chapter{Witt theory of quadratic spaces}
\label{ch:la01-witt}

Throughout this chapter \(k\) is a field with
\(\operatorname{char} k\ne 2\), and the ambient quadratic spaces are
finite dimensional and nondegenerate. For a quadratic form \(q\) we
use the symmetric bilinear form
\[
  B(x,y)=\frac{q(x+y)-q(x)-q(y)}{2},
  \qquad q(x)=B(x,x).
\]
Nondegeneracy means that \(B\) has zero radical. An isometry is a
linear isomorphism preserving \(q\), equivalently \(B\). We write
\(\perp\) for an orthogonal direct sum and \(\langle a\rangle\) for
the line with form \(x\mapsto ax^2\). These conventions agree with
the polarization and diagonalization of the archived notes
\cite[11.polarization and 11.diagonal\_congruence]{d05}.

For a subspace \(U\subseteq V\), write
\(U^\perp=\{x\in V:B(x,U)=0\}\).
When the restriction to \(U\) is nondegenerate, the decomposition
\(V=U\perp U^\perp\) is the symmetric case of the projection argument
in \cite[11.nondegenerate]{d05}. In fact, for \(v\in V\), the
restriction \(B(v,-)|_U\) has a unique representative \(u\in U\)
under \(U\cong U^*\) via \(B|_U\). Thus \(v-u\in U^\perp\).
The restriction to \(U^\perp\) is nondegenerate because a vector
orthogonal to both summands belongs to the radical of \(V\).

\section{Isotropic vectors and hyperbolic planes}

A nonzero vector \(e\) is \emph{isotropic} if \(q(e)=0\).
A space is \emph{anisotropic} when its only vector of value zero is
the zero vector, so the zero space is anisotropic. A subspace \(T\)
is \emph{totally isotropic} if \(q|_T=0\); by polarization this is
equivalent to \(B|_{T\times T}=0\). The hyperbolic plane \(H\) is
\(k^2\) with \(q_H(x,y)=xy\), equivalently
\(\langle 1\rangle\perp\langle -1\rangle\) by the invertible
change \((s,t)\mapsto(s+t,s-t)\).

\begin{lemma}[LA01-C06.01. Splitting a hyperbolic plane]
\label{lem:la01-hyperbolic-splitting}
If \(e\in V\) is isotropic, there is an isotropic vector \(f\)
with \(2B(e,f)=1\). Then \(P=ke+kf\) is isometric to \(H\) and
\(V=P\perp P^\perp\) with \(P^\perp\) nondegenerate.
\end{lemma}

\begin{proof}
Nondegeneracy and \(e\ne 0\) give a vector pairing nontrivially
with \(e\). Scale it to \(f_0\) with \(2B(e,f_0)=1\), and put
\(f=f_0-q(f_0)e\). The identities
\[
  2B(e,f)=1,\qquad
  q(f)=q(f_0)-2q(f_0)B(e,f_0)+q(f_0)^2q(e)=0
\]
show that \(e,f\) are independent and \(q(se+tf)=st\).
Their Gram matrix has diagonal zero and off-diagonal entries
\(1/2\), so the restriction to \(P\) is nondegenerate. Apply the
orthogonal splitting for a nondegenerate restriction.
\end{proof}

\section{Reflections and cancellation}

\begin{lemma}[LA01-C06.02. Moving vectors of the same nonzero value]
\label{lem:la01-reflections}
For \(w\in V\) with \(q(w)\ne 0\), the reflection
\[
  s_w(x)=x-\frac{2B(x,w)}{q(w)}\,w
\]
is an isometry. If \(u,v\in V\) satisfy
\(q(u)=q(v)=a\ne 0\), an isometry sending \(u\) to \(v\)
can be chosen as a product of at most two reflections.
\end{lemma}

\begin{proof}
The orthogonal decomposition \(V=kw\perp(kw)^\perp\) shows
that \(s_w\) negates the \(kw\) coordinate and fixes
\((kw)^\perp\). It is therefore an involution and preserves \(q\).
If \(q(u-v)\ne 0\), then \(2B(u,u-v)=q(u-v)\), so
\(s_{u-v}(u)=v\). Otherwise
\[
  q(u+v)=2q(u)+2q(v)-q(u-v)=4a\ne 0.
\]
In this case \(2B(u,u+v)=q(u+v)\), so \(s_{u+v}(u)=-v\).
The reflection \(s_v\) sends \(-v\) to \(v\). This second case
also works when \(u=v\).
\end{proof}

\begin{theorem}[LA01-C06.03. Witt cancellation]
\label{thm:la01-witt-cancellation}
Let \(U,V,W\) be nondegenerate quadratic spaces. If
\(U\perp W\cong V\perp W\), then \(U\cong V\).
The same conclusion holds when the copies of \(W\) are
replaced by two isometric nondegenerate spaces.
\end{theorem}

\begin{proof}
First let \(W=ke\) with \(q(e)=a\ne 0\). For an isometry
\(F:U\perp W\to V\perp W\), the vectors \(F(e)\) and the
specified \(e\) in the target both have value \(a\).
By Lemma~\ref{lem:la01-reflections} there is a target isometry
\(s\) sending \(F(e)\) to the target \(e\). The isometry \(sF\)
then carries the orthogonal complement of the source line to
the orthogonal complement of the target line. These
complements are \(U\) and \(V\).

For arbitrary \(W\), take an orthogonal basis by
\cite[11.diagonal\_congruence]{d05}. Nondegeneracy makes every
diagonal coefficient nonzero. Cancel the lines one at a time
by the preceding argument. When \(W=0\), the given isometry is
already the desired isometry. For two isometric copies of
\(W\), identify them by their given isometry before cancelling.
\end{proof}

\section{Witt decomposition and Witt index}

\begin{theorem}[LA01-C06.04. Witt decomposition and uniqueness]
\label{thm:la01-witt-decomposition}
There is an isometry
\[
  V\cong H^{\perp r}\perp A
\]
for an integer \(r\ge 0\) and an anisotropic space \(A\). The
integer \(r\) and the isometry class of \(A\) depend only on \(V\).
We call \(A\) the anisotropic part and \(r\) the Witt index.
Here \(H^{\perp 0}=0\).
\end{theorem}

\begin{proof}
If \(V\) is anisotropic, take \(r=0\). Otherwise split off a
hyperbolic plane by Lemma~\ref{lem:la01-hyperbolic-splitting} and
repeat in its nondegenerate complement. The dimension drops by
two at each step, so this process ends with an anisotropic
remainder, possibly zero.

For another decomposition \(H^{\perp s}\perp A'\), assume
\(r\le s\) after exchanging the two decompositions if needed.
Witt cancellation (Theorem~\ref{thm:la01-witt-cancellation})
gives \(A\cong H^{\perp(s-r)}\perp A'\). If \(s-r>0\), the
right side has a nonzero isotropic vector, contradicting
anisotropy of \(A\). Thus \(r=s\) and \(A\cong A'\).
\end{proof}

\begin{proposition}[LA01-C06.05. Totally isotropic subspaces and the Witt index]
\label{prop:la01-witt-index}
Let \(T\subseteq V\) be totally isotropic of dimension \(m\).
There is an orthogonal sum of \(m\) hyperbolic planes in \(V\)
whose chosen first isotropic basis vectors span \(T\).
Consequently \(m\le r\), and \(r\) is the maximum dimension of a
totally isotropic subspace. Every totally isotropic subspace
maximal by inclusion has dimension \(r\).
\end{proposition}

\begin{proof}
For \(m=0\) take the empty sum. For \(m>0\), choose
\(e\in T\setminus\{0\}\) and use
Lemma~\ref{lem:la01-hyperbolic-splitting} to write
\(V=P\perp P^\perp\) with \(P=ke+kf\) and \(2B(e,f)=1\).
For \(t\in T\), total isotropy gives \(B(t,e)=0\). The vector
\[
  t-2B(t,f)e
\]
also pairs to zero with \(f\), so it lies in \(T\cap P^\perp\).
Thus \(T=ke\oplus(T\cap P^\perp)\). Apply the induction
hypothesis to this \((m-1)\)-dimensional subspace inside the
nondegenerate space \(P^\perp\). This produces the claimed \(m\)
planes, with a nondegenerate orthogonal complement \(C\).

Decompose \(C\) by Theorem~\ref{thm:la01-witt-decomposition}.
Uniqueness gives \(r=m\) plus the Witt index of \(C\), hence
\(m\le r\). The span of one isotropic coordinate vector in each
of the \(r\) planes in a Witt decomposition has dimension \(r\).
If \(T\) is maximal by inclusion, \(C\) is anisotropic, since an
isotropic line in \(C\) could be adjoined orthogonally to \(T\).
Then the Witt index of \(C\) is zero and \(m=r\).
\end{proof}

For the general theory, see \cite[Ch.~II, Secs.~7--8]{witt-source}.

\begin{example}[LA01-C06.06. A Witt index that increases after extending scalars]
\label{ex:la01-witt-base-change}
On \(\mathbb{Q}^4\) take \(q(x,y,u,v)=xy+u^2-2v^2\).
The last two coordinates form a nondegenerate plane with no
nonzero isotropic vector. Indeed, \(u^2=2v^2\) with \(v\ne 0\)
would make \(2\) a rational square: writing \(u/v=a/b\) in
lowest terms would give \(a^2=2b^2\), forcing both \(a\) and
\(b\) even. If \(v=0\), then \(u=0\). The displayed form is
therefore \(H\) plus an anisotropic plane, and its Witt index
over \(\mathbb{Q}\) is \(1\).

Over \(\mathbb{Q}(\sqrt{2})\) the change
\((u,v)\mapsto(u+\sqrt{2}v,u-\sqrt{2}v)\) is invertible and
turns \(u^2-2v^2\) into a product of the two new coordinates.
That plane becomes hyperbolic, so the Witt index after extension
is \(2\). Both computations use
Theorem~\ref{thm:la01-witt-decomposition}.
\end{example}

\backmatter
{\raggedright
\begin{thebibliography}{9}

\bibitem{d05}
MathGenome, \emph{Linear Algebra and Foundational Higher Algebra}
(D05), revised notes archived in MathGenome v0.22.0.
Statement keys start with \texttt{K05}.

\bibitem{p103}
MathGenome, \emph{Algebra Core} (P1-03), Chapters~1--2, archived
in MathGenome v0.22.0.

\bibitem{d07}
MathGenome, \emph{Abstract Algebra} (D07), revised notes,
Chapter~16, archived in MathGenome v0.22.0.

\bibitem{p206}
MathGenome, \emph{Foundations of Numerical Analysis} (P2-06),
Chapter~5, archived in MathGenome v0.22.0.

\bibitem{d16}
MathGenome, \emph{Functional Analysis} (D16), the Hilbert-space
polar decomposition, archived in MathGenome v0.22.0.

\bibitem{commalg}
MathGenome, \emph{Advanced Commutative Algebra}, Results RK-03
and RK-04, archived in MathGenome v0.22.0.

\bibitem{stacks-sym}
The Stacks Project Authors,
\href{https://stacks.math.columbia.edu/tag/00DM}{\emph{Tensor algebra}},
Section~10.13, and
\href{https://stacks.math.columbia.edu/tag/01CI}{\emph{Symmetric and exterior powers}},
Lemma~17.21.3.

\bibitem{conrad-sym}
B. Conrad, \href{https://math.stanford.edu/~conrad/210APage/handouts/symmwedgebasis.pdf}{\emph{Bases of Symmetric and Exterior Powers}},
Math~210A, Stanford University, Theorem~1.2 and its proof.

\bibitem{witt-source}
R. Elman, N. Karpenko, and A. Merkurjev,
\href{https://www.math.ucla.edu/~rse/book/Kniga.pdf}{\emph{The Algebraic and Geometric Theory of Quadratic Forms}},
AMS Colloquium Publications 56, 2008, Chapter~II, Sections~7--8.

\bibitem{lam-witt}
T. Y. Lam,
\href{https://doi.org/10.1090/gsm/067}{\emph{Introduction to Quadratic Forms over Fields}},
Graduate Studies in Mathematics 67, AMS, 2005, Chapter~I, Section~4.

\end{thebibliography}
}
\end{document}
